% Options for packages loaded elsewhere
% Options for packages loaded elsewhere
\PassOptionsToPackage{unicode}{hyperref}
\PassOptionsToPackage{hyphens}{url}
\PassOptionsToPackage{dvipsnames,svgnames,x11names}{xcolor}
%
\documentclass[
  letterpaper,
  DIV=11,
  numbers=noendperiod]{scrartcl}
\usepackage{xcolor}
\usepackage{amsmath,amssymb}
\setcounter{secnumdepth}{5}
\usepackage{iftex}
\ifPDFTeX
  \usepackage[T1]{fontenc}
  \usepackage[utf8]{inputenc}
  \usepackage{textcomp} % provide euro and other symbols
\else % if luatex or xetex
  \usepackage{unicode-math} % this also loads fontspec
  \defaultfontfeatures{Scale=MatchLowercase}
  \defaultfontfeatures[\rmfamily]{Ligatures=TeX,Scale=1}
\fi
\usepackage{lmodern}
\ifPDFTeX\else
  % xetex/luatex font selection
\fi
% Use upquote if available, for straight quotes in verbatim environments
\IfFileExists{upquote.sty}{\usepackage{upquote}}{}
\IfFileExists{microtype.sty}{% use microtype if available
  \usepackage[]{microtype}
  \UseMicrotypeSet[protrusion]{basicmath} % disable protrusion for tt fonts
}{}
\makeatletter
\@ifundefined{KOMAClassName}{% if non-KOMA class
  \IfFileExists{parskip.sty}{%
    \usepackage{parskip}
  }{% else
    \setlength{\parindent}{0pt}
    \setlength{\parskip}{6pt plus 2pt minus 1pt}}
}{% if KOMA class
  \KOMAoptions{parskip=half}}
\makeatother
% Make \paragraph and \subparagraph free-standing
\makeatletter
\ifx\paragraph\undefined\else
  \let\oldparagraph\paragraph
  \renewcommand{\paragraph}{
    \@ifstar
      \xxxParagraphStar
      \xxxParagraphNoStar
  }
  \newcommand{\xxxParagraphStar}[1]{\oldparagraph*{#1}\mbox{}}
  \newcommand{\xxxParagraphNoStar}[1]{\oldparagraph{#1}\mbox{}}
\fi
\ifx\subparagraph\undefined\else
  \let\oldsubparagraph\subparagraph
  \renewcommand{\subparagraph}{
    \@ifstar
      \xxxSubParagraphStar
      \xxxSubParagraphNoStar
  }
  \newcommand{\xxxSubParagraphStar}[1]{\oldsubparagraph*{#1}\mbox{}}
  \newcommand{\xxxSubParagraphNoStar}[1]{\oldsubparagraph{#1}\mbox{}}
\fi
\makeatother


\usepackage{longtable,booktabs,array}
\newcounter{none} % for unnumbered tables
\usepackage{calc} % for calculating minipage widths
% Correct order of tables after \paragraph or \subparagraph
\usepackage{etoolbox}
\makeatletter
\patchcmd\longtable{\par}{\if@noskipsec\mbox{}\fi\par}{}{}
\makeatother
% Allow footnotes in longtable head/foot
\IfFileExists{footnotehyper.sty}{\usepackage{footnotehyper}}{\usepackage{footnote}}
\makesavenoteenv{longtable}
\usepackage{graphicx}
\makeatletter
\newsavebox\pandoc@box
\newcommand*\pandocbounded[1]{% scales image to fit in text height/width
  \sbox\pandoc@box{#1}%
  \Gscale@div\@tempa{\textheight}{\dimexpr\ht\pandoc@box+\dp\pandoc@box\relax}%
  \Gscale@div\@tempb{\linewidth}{\wd\pandoc@box}%
  \ifdim\@tempb\p@<\@tempa\p@\let\@tempa\@tempb\fi% select the smaller of both
  \ifdim\@tempa\p@<\p@\scalebox{\@tempa}{\usebox\pandoc@box}%
  \else\usebox{\pandoc@box}%
  \fi%
}
% Set default figure placement to htbp
\def\fps@figure{htbp}
\makeatother





\setlength{\emergencystretch}{3em} % prevent overfull lines

\providecommand{\tightlist}{%
  \setlength{\itemsep}{0pt}\setlength{\parskip}{0pt}}



 


\KOMAoption{captions}{tableheading}
\makeatletter
\@ifpackageloaded{caption}{}{\usepackage{caption}}
\AtBeginDocument{%
\ifdefined\contentsname
  \renewcommand*\contentsname{Table of contents}
\else
  \newcommand\contentsname{Table of contents}
\fi
\ifdefined\listfigurename
  \renewcommand*\listfigurename{List of Figures}
\else
  \newcommand\listfigurename{List of Figures}
\fi
\ifdefined\listtablename
  \renewcommand*\listtablename{List of Tables}
\else
  \newcommand\listtablename{List of Tables}
\fi
\ifdefined\figurename
  \renewcommand*\figurename{Figure}
\else
  \newcommand\figurename{Figure}
\fi
\ifdefined\tablename
  \renewcommand*\tablename{Table}
\else
  \newcommand\tablename{Table}
\fi
}
\@ifpackageloaded{float}{}{\usepackage{float}}
\floatstyle{ruled}
\@ifundefined{c@chapter}{\newfloat{codelisting}{h}{lop}}{\newfloat{codelisting}{h}{lop}[chapter]}
\floatname{codelisting}{Listing}
\newcommand*\listoflistings{\listof{codelisting}{List of Listings}}
\makeatother
\makeatletter
\makeatother
\makeatletter
\@ifpackageloaded{caption}{}{\usepackage{caption}}
\@ifpackageloaded{subcaption}{}{\usepackage{subcaption}}
\makeatother
\usepackage{bookmark}
\IfFileExists{xurl.sty}{\usepackage{xurl}}{} % add URL line breaks if available
\urlstyle{same}
\hypersetup{
  pdftitle={Executive Summary --- Cell-Type-Specific Biomarkers in Human Liver Fibrosis},
  colorlinks=true,
  linkcolor={blue},
  filecolor={Maroon},
  citecolor={Blue},
  urlcolor={Blue},
  pdfcreator={LaTeX via pandoc}}


\title{Executive Summary --- Cell-Type-Specific Biomarkers in Human
Liver Fibrosis}
\author{}
\date{2026-06-01}
\begin{document}
\maketitle

\renewcommand*\contentsname{Table of contents}
{
\setcounter{tocdepth}{3}
\tableofcontents
}

\section{Objective}\label{objective}

Discover and prioritize \textbf{cell-type-specific biomarkers and
therapeutic targets} in human liver fibrosis from single-cell
transcriptomics (primary: GSE136103, cirrhotic vs healthy; validation:
GSE244832, MASLD/MASH F0--F4), focusing on the hepatic
stellate/myofibroblast, macrophage/monocyte, and endothelial
compartments.

\section{Approach (one paragraph)}\label{approach-one-paragraph}

A reproducible Snakemake + Scanpy pipeline ingests and QCs the data,
integrates over donor (not disease) and verifies fibrosis signal is
preserved, annotates and validates the three disease compartments, runs
\textbf{donor-aware pseudobulk} differential expression, characterizes
mechanism via pathway/TF activity and ligand--receptor analysis (with
downstream-target corroboration), and prioritizes candidates with a
transparent composite score combining DE effect, cell-type specificity,
cross-dataset reproducibility, druggability, biomarker accessibility,
and an explainable ML (XGBoost + SHAP) signal.

\section{Top candidates}\label{top-candidates}

{\def\LTcaptype{none} % do not increment counter
\begin{longtable}[]{@{}llllll@{}}
\toprule\noalign{}
& rank & gene & compartment & composite & direction \\
\midrule\noalign{}
\endhead
\bottomrule\noalign{}
\endlastfoot
0 & 1.0 & PDGFRA & stellate\_myofibroblast & 0.671195 & up \\
1 & 2.0 & ACKR1 & endothelial & 0.634854 & up \\
2 & 3.0 & LUM & stellate\_myofibroblast & 0.624039 & up \\
3 & 4.0 & DCN & stellate\_myofibroblast & 0.611797 & up \\
4 & 5.0 & TREM2 & macrophage\_monocyte & 0.606202 & up \\
5 & 6.0 & PLVAP & endothelial & 0.580863 & up \\
7 & 7.0 & COL3A1 & stellate\_myofibroblast & 0.552956 & up \\
9 & 8.0 & MMP7 & cholangiocyte & 0.545108 & up \\
10 & 9.0 & SPP2 & cholangiocyte & 0.543718 & up \\
11 & 10.0 & MMP2 & stellate\_myofibroblast & 0.542169 & up \\
12 & 11.0 & COL1A1 & stellate\_myofibroblast & 0.537588 & up \\
15 & 12.0 & RGCC & endothelial & 0.511566 & up \\
16 & 13.0 & COL1A2 & stellate\_myofibroblast & 0.505817 & up \\
17 & 14.0 & RGS4 & cholangiocyte & 0.494848 & up \\
22 & 15.0 & VWF & endothelial & 0.477451 & up \\
33 & 16.0 & FABP4 & endothelial & 0.431836 & up \\
34 & 17.0 & RGS10 & macrophage\_monocyte & 0.428405 & up \\
49 & 18.0 & IER3 & macrophage\_monocyte & 0.405677 & up \\
78 & 19.0 & PLXDC2 & macrophage\_monocyte & 0.390093 & up \\
118 & 20.0 & CTSH & macrophage\_monocyte & 0.362371 & up \\
\end{longtable}
}

\section{Translational
interpretation}\label{translational-interpretation}

The ranked table flags, per compartment, which candidates are most
relevant for \textbf{diagnostics} (secreted/accessible, reproducible
across datasets) versus \textbf{therapeutic targeting} (druggable,
mechanistically central in the fibrotic crosstalk). See the full report
and \texttt{results/08\_score/candidate\_scores.tsv} for the
per-candidate rationale.




\end{document}
